% 2007 年四川卷文科第 19 题：题图。
% 来源 PDF：普通高考/2007/2007四川文.pdf，SHA256 1318df44733ec939e8789379d6fa2115a2c6a63e84b73c1f2bf3c080007ce68c
% 原页：第 2 页（A3 三栏）第 1 栏，印刷页 412，图在题干下方。
% 重绘依据：P、M 在上方且 PM 水平，C 在 P 正下方，B 在 C 右侧，A 在左下方；可见线段实线，被梯形遮住的 PC、CB、AC 虚线，C 处画直角记号；与文科原页逐线比对一致。
\begin{tikzpicture}[scale=2.2, semithick,
  every node/.style={inner sep=1.5pt}]
  \coordinate (P) at (0,3);
  \coordinate (M) at (1,3);
  \coordinate (C) at (0,2);
  \coordinate (B) at (2,2);
  \coordinate (A) at (-0.95,1.55);
  % 被遮挡的三条线
  \draw[dashed] (P) -- (C);
  \draw[dashed] (C) -- (B);
  \draw[dashed] (A) -- (C);
  % 可见线段
  \draw (P) -- (M);
  \draw (P) -- (A);
  \draw (M) -- (B);
  \draw (A) -- (B);
  \draw (A) -- (M);
  % C 处直角记号
  \draw ($(C)+(0.18,0)$) -- ($(C)+(0.18,0.18)$) -- ($(C)+(0,0.18)$);
  \node[anchor=south east] at ([xshift=-1pt,yshift=1pt]P) {$P$};
  \node[anchor=south west] at ([xshift=1pt,yshift=1pt]M) {$M$};
  \node[anchor=north west] at ([xshift=2pt,yshift=-2pt]C) {$C$};
  \node[anchor=north west] at ([xshift=1pt,yshift=-1pt]B) {$B$};
  \node[anchor=north east] at ([xshift=-1pt,yshift=-1pt]A) {$A$};
\end{tikzpicture}
